% =========================
%  Insulation Coordination
% =========================
\intermezzo{Lightning}{}

\begin{frame}{Lightning stroke}
  \begin{columns}[c,onlytextwidth]
    \begin{column}{0.48\textwidth}
      \centering
      \topoimg[width=\linewidth,height=0.66\textheight,keepaspectratio]{\icfig{m08_image5}}
      {\color{atGray}\footnotesize IEEE Std 998-1996}
    \end{column}

    \begin{column}{0.48\textwidth}
      \begin{itemize}
        \setlength{\itemsep}{0.35em}
        \item When a lightning strikes an overhead line, the current wave propagates in two directions
        \item The travelling current wave creates also a voltage wave travelling along with the current wave
        \item The voltage wave depends on the lightning current and the surge impedance of the line
        \item Surge impedance is quite different from DC or 50/60Hz impedances, is not associated with any frequency or line length
      \end{itemize}

    \end{column}
  \end{columns}
  \vspace{0.3em}
\end{frame}

\begin{frame}{Lightning stroke}
  \begin{itemize}
    \item Surge impedance depends on the line configuration. Typical values are:\\
    \quad OHL:\quad$400 - 600  \Omega$ \\
    \quad Cable:\quad$20 - 60  \Omega$
    \item The voltage wave travelling with the lightning stroke is defined by: \\
    \quad $E=\frac{1}{3} \cdot I_S \cdot Z_S$
    \quad

  \end{itemize}
\end{frame}

\begin{frame}{Travelling wave}
  \begin{columns}[c,onlytextwidth]
    \begin{column}{0.56\textwidth}
      \begin{itemize}
        \setlength{\itemsep}{0.4em}
        \item Travelling-wave propagation is not comparable to the current
              flow at nominal frequency.
        \item On an open circuit no current flows, but the travelling wave
              can still enter the conductor.
        \item When the wave hits the open circuit (e.g. a transformer) the
              surge is reflected: the voltage keeps its polarity but the
              amplitude doubles.
        \item The increased amplitude at transformers must be considered in
              the equipment design.
      \end{itemize}
      \vspace{0.2em}
    \end{column}
    \begin{column}{0.40\textwidth}
      \centering
      \topoimg[width=\linewidth,height=0.66\textheight,keepaspectratio]{\icfig{m08_image6}}
      {\color{atGray}\footnotesize Westinghouse, ``Electrical Transmission and Distribution Reference Book'', 1961}
    \end{column}
  \end{columns}
\end{frame}

\begin{frame}{Energy}
  \begin{columns}[c,onlytextwidth]
    \begin{column}{0.56\textwidth}
      \begin{itemize}
        \setlength{\itemsep}{0.35em}
        \item Energy is defined not only by the peak current but by the total
              charge $Q$ of the lightning stroke.
        \item For a $30 kA$ stroke and a line surge impedance of $600 \Omega$,
              the travelling-wave voltage is $\sim9000kV$.
        \item Overvoltage causes flashover on line insulators, diverting most
              of the charge to ground before it reaches protected equipment.
        \item The most critical strokes hit substations directly --- proper
              lightning protection of equipment must be ensured.
      \end{itemize}
    \end{column}
    \begin{column}{0.40\textwidth}
      \centering
      \topoimg[width=\linewidth,height=0.66\textheight,keepaspectratio]{\icfig{m08_image13}}
      $Q=\int i\,dt$\\
      $i2=\int i^2\,dt$
    \end{column}
  \end{columns}
\end{frame}

